\section*{Abstract}

Large Language Models (LLMs) are increasingly integrated into applications that
process external data and interact with tools, thereby expanding their attack
surface. One of the most important threats to such systems is prompt injection,
where malicious instructions embedded in untrusted content manipulate model
behaviour. Since current LLM architectures process instructions and data within
the same context window, distinguishing between trusted and untrusted
information remains a fundamental challenge.

The central question is how effective the separation of instructions and data
is as a defence strategy against prompt injection attacks. To address this
question, the work reviews the underlying threat model, structured query
approaches such as StruQ and SecAlign, and additional architectural and
system-level defence mechanisms. Current evaluation methodologies and recent
adaptive attacks are also analysed to assess the limits of existing approaches.

The results indicate that separating instructions and data significantly
improves robustness against many attacks, but does not eliminate the underlying
vulnerability. Fine-tuning-based approaches remain susceptible to adaptive
attacks, while stronger guarantees generally require architectural isolation
and defence-in-depth principles.

Consequently, prompt injection remains an open problem. Secure deployment of
LLM-integrated systems requires not only improved models, but also stronger
system architectures and more rigorous evaluation methodologies.